% fixture: babel/french-lists
% tier: T3
% engine: document
% covers: french.ldf's lists as the main language: \listFB's vertical settings (.8ex and .4ex, in the list's font: the abstract's \small too), itemize's \itemizeFB (no skips, less \parsep in a list's item, \textemdash labels, \leftmargin from \labelwidthFB, \labelsep and \listindentFB), enumerate and description with \FB@listHsettings, quote (\listparindent zero for a list in it), quotation, verse and thebibliography with \listFB
% issue: #338
\documentclass{article}
\usepackage[T1]{fontenc}
\usepackage[french]{babel}
\begin{document}
\begin{abstract}
Un résumé court qui parle de la théorie de l'information et de son application aux communications numériques modernes.
\end{abstract}
Texte avant les listes.
\begin{itemize}
\item Premier élément de la liste avec un texte assez long pour remplir plus d'une ligne entière de la page.
\item Deuxième élément.
\begin{itemize}
\item Imbriqué.
\begin{itemize}
\item Encore plus profond.
\end{itemize}
\end{itemize}
\end{itemize}
Puis une énumération :
\begin{enumerate}
\item Un.
\item Deux, avec une sous-liste :
\begin{enumerate}
\item alpha ;
\item beta.
\end{enumerate}
\item Trois, avec des puces :
\begin{itemize}
\item une ;
\item deux.
\end{itemize}
\end{enumerate}

\begin{description}
\item[Terme] Sa définition, qui est longue et qui continue sur la ligne suivante pour voir l'indentation.
\item[Autre] Encore.
\end{description}
\begin{quote}
Une citation courte.
\begin{itemize}
\item Un élément de deux paragraphes, assez long pour remplir la ligne entière de la citation.

Le second paragraphe de l'élément.
\end{itemize}
\end{quote}
\begin{quotation}
Une longue citation en plusieurs paragraphes.

Le second paragraphe.
\end{quotation}
\begin{verse}
Un vers\\
et un autre.
\end{verse}
\begin{thebibliography}{9}
\bibitem{a} Un ouvrage.
\end{thebibliography}
\end{document}
